\documentclass[10pt,a4paper]{amsart}
\usepackage[margin=19mm]{geometry}
\usepackage{amsmath,amssymb,mathtools}
\usepackage[scr=rsfs,cal=euler]{mathalfa}
\usepackage{xfrac}
\usepackage[unicode,hidelinks,bookmarks=false]{hyperref}
\newcommand{\ie}{i.e.\ }
\newcommand{\cf}{cf.\ }
\newcommand{\resp}{respectively\ }
\newcommand{\Subsection}{\subsection}
\providecommand{\nobreakdash}{\nobreak\hbox{-}}
\setlength{\emergencystretch}{2em}
\multlinegap0pt
\usepackage{amssymb}
\usepackage{xcolor}
\usepackage{caption}
\usepackage{float}
\usepackage{subfig}
\newtheorem{definition}{Definition}
\newtheorem{lemma}{Lemma}
\def\C {\mathcal{C}}
\title{Multi-term fractional linear equations\\ modeling oxygen subdiffusion through capillaries}
\author{Vittorino Pata, Sergii Siryk, Nataliya Vasylyeva}
\date{Source: arXiv:2210.05009v1 (CC BY 4.0)}
\begin{document}
\maketitle

\noindent\emph{Source and licence.} Multi-term fractional linear equations\\ modeling oxygen subdiffusion through capillaries. Version arXiv:2210.05009v1 (CC BY 4.0), distributed by arXiv under the Creative Commons Attribution 4.0 International licence, http://creativecommons.org/licenses/by/4.0/, as recorded in the arXiv licence metadata for this exact version and shown on the abstract page https://arxiv.org/abs/2210.05009v1. Full licence notice: \texttt{licenses/arxiv-2210.05009-CC-BY-4.0.txt}.\par\smallskip

\noindent\textit{Editorial scope.} Counted unit: the article's lemma l4.1 (source0.tex lines 618--654). The article's complete original proof of it is delivered with it (source0.tex lines 655--811, one \texttt{proof} environment of 157 lines). No mathematics is written, repaired or re-derived: the statement and the proof are the article's own lines, in the article's own notation, and no retained line is substituted or re-flowed.  Retained uncounted as the source-local context the proof needs: lines 288--303; lines 815--855.  Nothing was omitted from any retained span, so the package carries no omission record.  The retained lines cite 1 works, whose entries are transcribed verbatim from the article's own bibliography source in the e-print and delivered with short editorial labels recorded in the item's reference mapping.  No retained line contains a figure environment, an image-inclusion command or drawing code, so no image is delivered and the package declares no asset dependency.  Editorial formatting only, in addition: an AMS \texttt{amsart} preamble that restates the article's own declarations and macros used by the retained lines, namely the environments \texttt{definition}, \texttt{lemma} on separate counters, together with the standard packages those lines need.  The retained environments print in this document as Definition 1, Lemma 1 in order of appearance, whereas the article prints the same items with the numbering of its own full document.  The complete native source of the article as distributed in the arXiv e-print for this exact version is bundled unchanged as \texttt{sources/131868/source0.tex}.
\begin{definition}\label{d2.1}
 A function $v=v(x,t)$ belongs to the class $\C^{l+\alpha,\frac{l+\alpha}{2}\nu}(\bar{\Omega}_{T})$, for $l=0,1,2,$ if the function $v$ together with its corresponding derivatives are continuous and the norms here below are finite:
\begin{align*}
\|v\|_{\C^{l+\alpha,\frac{l+\alpha}{2}\nu}(\bar{\Omega}_{T})} &=
\|v\|_{\C([0,T],
\C^{l+\alpha}(\bar{\Omega}))}+\sum_{|j|=0}^{l}\langle
D_{x}^{j}v\rangle^{(\frac{l+\alpha-|j|}{2}\nu)}_{t,{\Omega}_{T}},
\quad
l=0,1,\\
\|v\|_{\C^{2+\alpha,\frac{2+\alpha}{2}\nu}(\bar{\Omega}_{T})}
&=\|v\|_{\C([0,T],
\C^{2+\alpha}(\bar{\Omega}))}+\|\mathbf{D}_{t}^{\nu}v\|_{\C^{\alpha,\frac{\alpha}{2}\nu}
(\bar{\Omega}_{T})}+\sum_{|j|=1}^{2}\langle
D_{x}^{j}v\rangle^{(\frac{2+\alpha-|j|}{2}\nu)}_{t,{\Omega}_{T}}.
\end{align*}
\end{definition}

\begin{lemma}\label{l4.1bis}
Let the assumptions of Lemma \ref{l4.1} hold. Besides,
let $\gamma>\frac{(3+\alpha)\nu}{2}$ and
\[
\mathbf{D}_{t}^{\nu}w_1\in\C_{0}^{1+\alpha,\frac{1+\alpha}{2}\nu}(\partial\Omega_{T}) \qquad
\text{and}\qquad w_{2}\in\C^{\gamma}([0,T],\C^{2}(\partial\Omega)).
\]
If $\gamma<1$ we
 also require $\mathbf{D}_{t}^{\nu}w_2\in\C_{0}^{1+\alpha,\frac{1+\alpha}{2}\nu}(\partial\Omega_{T})$.
 Then, for $\beta\in(0,\frac{\nu(1-\alpha)}{2})$, with $\delta$ as above,
 the following estimates hold:
 \begin{description}
	\item[i]
	\begin{align*}
	&\|w_2\xi \mathbf{D}_{t}^{\beta}w_1\|_{\C^{1+\alpha,\frac{1+\alpha}{2}\nu}(\partial\Omega^{\star}_{\tau})}\\&\leq C
		[\tau^{\nu-\beta}+\tau^{\nu\frac{(1-\alpha)}{2}-\beta}+\tau^{\nu-\beta}r^{-1-\alpha}+\tau^{\nu-\beta}r^{-\alpha}]\|\mathbf{D}_{t}^{\nu}w_1\|_{\C^{1+\alpha,\frac{1+\alpha}{2}\nu}(\partial\Omega^{\star}_{\tau})}.
	\end{align*}
		\item[ii] If $w_2(x,0)=0$ then
		\begin{align*}
		&\|w_2\xi \mathbf{D}_{t}^{\nu}w_1\|_{\C^{1+\alpha,\frac{1+\alpha}{2}\nu}(\partial\Omega^{\star}_{\tau})}\\&
		\leq C
		[\tau^{\delta-\frac{(\alpha+1)}{2}\nu}+\tau^{\delta}r^{-\alpha-1}+\tau^{\delta-\frac{\alpha\nu}{2}}r^{-1}][\|\mathbf{D}_{t}^{\nu}w_1\|_{\C^{1+\alpha,\frac{1+\alpha}{2}\nu}(\partial\Omega^{\star}_{\tau})}+\|w_1\|_{\C^{2+\alpha,\frac{2+\alpha}{2}\nu}(\bar{\Omega}^{r}_{\tau})}].
		\end{align*}
		\item[iii]
		\begin{align*}
		\|\xi\mathfrak{J}_{\beta}(t)\|_{\C^{1+\alpha,\frac{1+\alpha}{2}\nu}(\partial\Omega^{\star}_{\tau})}&\leq C
		\frac{\tau^{\delta-\beta}}{r^{1+\alpha}}[\tau^{\nu}+\tau^{\beta}+r^{\alpha}(\tau^{\beta-\nu\alpha/2}+\tau^{\nu\frac{(1+\alpha)}{2}}+\tau^{\nu\frac{2-\alpha}{2}}+r\tau^{\nu\frac{1-\alpha}{2}}]\\
		&
		\quad\cdot
		[\|\mathbf{D}_{t}^{\nu}w_1\|_{\C^{1+\alpha,\frac{1+\alpha}{2}\nu}(\partial\Omega^{\star}_{\tau})}+\|w_1\|_{\C^{2+\alpha,\frac{2+\alpha}{2}\nu}(\bar{\Omega}^{r}_{\tau})}],
		\end{align*}
				\begin{align*}
		\|\xi\mathfrak{J}_{\nu}(t)\|_{\C^{1+\alpha,\frac{1+\alpha}{2}\nu}(\partial\Omega^{\star}_{\tau})}&\leq C
		\frac{\tau^{\delta-\nu}}{r^{1+\alpha}}[\tau^{\nu}+r^{\alpha}(\tau^{\nu-\nu\alpha/2}+\tau^{\nu\frac{(1+\alpha)}{2}}+\tau^{\nu\frac{2-\alpha}{2}}+r\tau^{\nu\frac{1-\alpha}{2}}]\\
		&
		\quad\cdot
		[\|\mathbf{D}_{t}^{\nu}w_1\|_{\C^{1+\alpha,\frac{1+\alpha}{2}\nu}(\partial\Omega^{\star}_{\tau})}+\|w_1\|_{\C^{2+\alpha,\frac{2+\alpha}{2}\nu}(\bar{\Omega}^{r}_{\tau})}].
		\end{align*}
		\end{description}
Again, $C$ depends only on $\nu,\beta,T$, the Lebesgue measure of $\Omega$ and the norm of $w_2$.
\end{lemma}

\begin{lemma}\label{l4.1}
Let $x^{0}\in\bar{\Omega}$ be arbitrarily fixed,
let $0<\beta<\nu\leq 1.$
We assume that
$$w_1\in\C_{0}^{2+\alpha,\frac{2+\alpha}{2}\nu}(\bar{\Omega}_{T})
\qquad\text{and}\qquad w_2\in\C^{\gamma}([0,T],\C^{1}(\bar{\Omega})),$$
with $\gamma\geq\frac{2+\alpha}{2}\nu$, and we set
$$\delta=\min\{1,\gamma\}.$$
If $\gamma<1$ we additionally  require  $\mathbf{D}_{t}^{\nu}w_2\in\C^{\alpha,\alpha\nu/2}(\bar{\Omega}_{T})$.
Then, for
any $\beta\in(0,\frac{2-\alpha}{2}\nu)$ and any $\tau\in(0,T]$, the following estimates hold:
\begin{description}
	\item[i]
\begin{align*}
	&\|w_2\xi \mathbf{D}_{t}^{\beta}w_1\|_{\C^{\alpha,\alpha\nu/2}(\bar{\Omega}^{r}_{\tau})}\\
	&\leq C[\tau^{\nu-\beta+\alpha\nu/2}+\tau^{\nu-\beta+\nu\alpha/2}r^{-\alpha}+\tau^{\nu-\beta}+\tau^{\nu-\beta}r^{-\alpha}]\|\mathbf{D}_{t}^{\nu}w_1\|_{\C^{\alpha,\alpha\nu/2}(\bar{\Omega}^{r}_{\tau})}.
\end{align*}
		\item[ii] If $w_2(x,0)=0$ then
		\[
		\|w_2\xi \mathbf{D}_{t}^{\nu}w_1\|_{\C^{\alpha,\alpha\nu/2}(\bar{\Omega}^{r}_{\tau})}\leq C[\tau^{\delta-\alpha\nu/2}+\tau^{\delta}r^{-\alpha}+\tau^{\delta}]\|\mathbf{D}_{t}^{\nu}w_1\|_{\C^{\alpha,\alpha\nu/2}(\bar{\Omega}^{r}_{\tau})}.
		\]
		\item[iii]
		\[
		\|\xi\mathfrak{J}_{\beta}(t)\|_{\C^{\alpha,\alpha\nu/2}(\bar{\Omega}^{r}_{\tau})}\leq C
		\tau^{\delta-\beta+\nu(1-\alpha)/2}[\tau^{\nu/2}+\tau^{\nu\alpha}+\tau^{\nu\frac{(1+\alpha)}{2}}(r^{-\alpha}+1)]
		\|w_1\|_{\C^{2+\alpha,\frac{2+\alpha}{2}\nu}(\bar{\Omega}^{r}_{\tau})},
		\]
		and
				\[
		\|\xi\mathfrak{J}_{\nu}(t)\|_{\C^{\alpha,\alpha\nu/2}(\bar{\Omega}^{r}_{\tau})}\leq C
\tau^{\delta-\nu(1+\alpha)/2}[\tau^{\nu/2}+\tau^{\nu\alpha}+\tau^{\nu\frac{(1+\alpha)}{2}}(r^{-\alpha}+1)]
		\|w_1\|_{\C^{2+\alpha,\frac{2+\alpha}{2}\nu}(\bar{\Omega}^{r}_{\tau})}.
		\]
		\end{description}
			
The positive quantity $C$ depends only on $\nu,\beta,T$,  the Lebesgue measure of $\Omega$ and the norm of $w_2$.
\end{lemma}

\begin{proof}
%The proof of Lemma \ref{l4.1bis} is very similar to the one of Lemma \ref{l4.1}.
%Accordingly, we restrict ourselves to the verification of points (i)-(iii).
We start with the evaluation of the term $\mathbf{D}_{t}^{\beta}w_1$. To this end, appealing to representation (10.34) in \cite{KPSV1}, we deduce that
\[
w_2\xi\mathbf{D}_{t}^{\beta}w_1=i_1+i_2,
\]
where we put
\begin{align*}
i_1&=i_1(x,t)=\xi(x)w_{2}(x,t)\frac{{t}^{\nu-\beta}}{\Gamma(1+
\nu-\beta)}\mathbf{D}_{t}^{\nu} w_1(x,t),\\
i_2&=i_2(x,t)=\xi(x)w_{2}(x,t)\int_{0}^{t}\frac{(t-s)^{\nu-\beta-1}}{\Gamma(\nu-\beta)}[\mathbf{D}_{t}^{\nu}w_{1}(x,t)-\mathbf{D}_{s}^{\nu}w_{1}(x,s)]ds.
\end{align*}
We estimate the norms of  $i_1$ and $i_2$ separately.

\smallskip
\noindent
As for $\|i_1\|_{\C^{\alpha,\alpha\nu/2}(\bar{\Omega}^{r}_{\tau})},$ the desired bound is a simple consequence of the following easily verified relations:
\begin{align*}
\underset{\bar{\Omega}^{r}_{\tau}}{\sup}\, |i_1|&\leq C\tau^{\nu-\beta}\underset{\bar{\Omega}^{r}_{\tau}}{\sup}\,|w_2|\,
\underset{\bar{\Omega}^{r}_{\tau}}{\sup}\,|\mathbf{D}_{t}^{\nu}w_{1}|,\\
\langle i_1\rangle_{x,\Omega^{r}_{\tau}}^{(\alpha)}&\leq\tau^{\nu-\beta}\Big[ \langle w_2\rangle_{x,\bar{\Omega}^{r}_{\tau}}^{(\alpha)}
\underset{\bar{\Omega}^{r}_{\tau}}{\sup}\,|\mathbf{D}_{t}^{\nu}w_{1}|
+\underset{\bar{\Omega}^{r}_{\tau}}{\sup}\,|w_2|
\langle \xi \mathbf{D}_{t}^{\nu}w_{1}\rangle_{x,\bar{\Omega}^{r}_{\tau}}^{(\alpha)}\Big]\\
&
\leq C[r^{1-\alpha}\tau^{\nu-\beta}+\tau^{\nu-\beta}r^{-\alpha}+\tau^{\nu-\beta}]\|\mathbf{D}_{t}^{\nu}w_{1}\|_{\C([0,\tau],\C^{\alpha}(\bar{\Omega}))},
\end{align*}
and
\begin{align*}
\langle i_1\rangle_{t,\Omega^{r}_{\tau}}^{(\alpha\nu/2)}&\leq C\Big[\tau^{\nu-\beta+\delta-\alpha\nu/2}\langle w_2\rangle_{t,\Omega^{r}_{\tau}}^{(\delta)}\|\mathbf{D}_{t}^{\nu}w_1\|_{\C(\bar{\Omega}_{\tau}^{r})}
+
\langle t^{\nu-\beta} \mathbf{D}_{t}^{\nu}w_1\rangle_{t,\Omega^{r}_{\tau}}^{(\alpha\nu/2)}\|w_2\|_{\C(\bar{\Omega}_{T})}\Big]\\
&
\leq
C\tau^{\nu-\beta-\alpha\nu/2}[1+\tau^{\alpha\nu/2}+\tau^{\delta}]\|w_2\|_{\C^{\gamma}([0,T],\C^{1}(\bar{\Omega}))}
\|\mathbf{D}_{t}^{\nu}w_1\|_{\C^{\alpha\nu/2}([0,\tau],\C(\bar{\Omega}^{r}))}.
\end{align*}

\noindent
Coming to $\|i_2\|_{\C^{\alpha,\alpha\nu/2}(\bar{\Omega}_{\tau}^{r})}$, we easily find
\[
\|i_2\|_{\C(\bar{\Omega}_{\tau}^{r})}+\langle i_2\rangle_{x,\Omega_{\tau}^{r}}^{(\alpha)}
\leq C\tau^{\nu-\beta}\Big[\tau^{\alpha\nu/2}(1+r^{-\alpha})\langle\mathbf{D}_{t}^{\nu}w_1\rangle_{t,\Omega_{\tau}^{r}}^{(\alpha\nu/2)}
+
\langle\mathbf{D}_{t}^{\nu}w_1\rangle_{x,\Omega_{\tau}^{r}}^{(\alpha)}
\Big]\|w_2\|_{\C([0,T],\C^{\alpha}(\bar{\Omega}))}.
\]
In order to complete the estimate of $i_2$, hence establishing
point (i), we are left to examine the difference $|i_2(x,t_2)-i_2(x,t_1)|$. To this end, assuming $t_2>t_1$ and setting $\Delta t=t_2-t_1$, we have
\[
|i_2(x,t_2)-i_2(x,t_1)|\leq \sum_{j=1}^{4}i_{2,j},
\]
where
\begin{align*}
i_{2,1}&=\xi|w_{2}(x,t_2)-w_{2}(x,t_1)|\Big|\int_{0}^{t_2}\frac{s^{\nu-\beta-1}}{\Gamma(\nu-\beta)}[\mathbf{D}_{t}^{\nu}w_1(x,t_{2}-s)-
\mathbf{D}_{t}^{\nu}w_1(x,t_{2})]ds\Big|,\\
i_{2,2}&=\xi|w_{2}(x,t_1)|\Big|\int_{0}^{t_1}\frac{s^{\nu-\beta-1}}{\Gamma(\nu-\beta)}[\mathbf{D}_{t}^{\nu}w_1(x,t_{2}-s)-
\mathbf{D}_{t}^{\nu}w_1(x,t_{1}-s)]ds\Big|,\\
i_{2,3}&=\xi|w_{2}(x,t_1)||\mathbf{D}_{t}^{\nu}w_1(x,t_{2})-\mathbf{D}_{t}^{\nu}w_1(x,t_{1})|\bigg|\int_{0}^{t_1}\frac{s^{\nu-\beta-1}}{\Gamma(\nu-\beta)}ds\bigg|,\\
i_{2,4}&=\xi|w_{2}(x,t_1)|\bigg|\int_{t_1}^{t_2}\frac{s^{\nu-\beta-1}}{\Gamma(\nu-\beta)}[\mathbf{D}_{t}^{\nu}w_1(x,t_{2}-s)-\mathbf{D}_{t}^{\nu}w_1(x,t_{2})]ds\bigg|.
\end{align*}
Exploiting the smoothness of the functions $w_2$ and $\mathbf{D}_{t}^{\nu}w_1$, and taking into account of the relation between $\nu$ and $\beta$, we arrive at the sought estimate
for $\langle i_2\rangle_{t,\Omega_{\tau}^{r}}^{(\alpha\nu/2)}$.

To verify statement (ii), it is worth noting that
the assumptions on $w_2$ (i.e. $w_2(x,0)=0$) provide the estimate
\[
\|w_2\|_{\C^{\alpha,\alpha\nu/2}(\bar{\Omega}_{\tau}^{r})}\leq C\tau^{\delta-\alpha\nu/2}
[\tau^{\alpha\nu/2}+1+r^{1-\alpha}\tau^{\alpha\nu/2}]\|w_2\|_{\C^{\delta}([0,T],\C^{1}(\bar{\Omega}))}.
\]
Collecting this bound with the regularity of $\mathbf{D}_{t}^{\nu}w_1$, the desired claim follows.

Coming to (iii), we restrict ourselves to the verification of the first inequality, for the second one is deduced in a similar manner.
Straightforward calculations lead to the relations
\[
\underset{\bar{\Omega}_{\tau}^{r}}{\sup}\,|\xi\mathfrak{J}_{\beta}(t)|\leq C\tau^{\delta-\beta+\nu}
\langle w_2\rangle_{t,\Omega_{\tau}^{r}}^{(\delta)}\|\mathbf{D}_{t}^{\nu}w_1\|_{\C(\bar{\Omega}_{\tau}^{r})},
\]
\begin{align*}
&\Big|\xi(x_2)\mathfrak{J}_{\beta}(t)|_{x=x_1}-\xi(x_1)\mathfrak{J}_{\beta}(t)|_{x=x_2}\Big|
\\
&
\leq C r^{-\alpha}\tau^{\delta-\beta+\nu}
\|\mathbf{D}_{t}^{\nu}w_1\|_{\C(\bar{\Omega}_{\tau}^{r})}
\langle w_2\rangle_{t,\Omega_{\tau}^{r}}^{(\delta)}|x_1-x_2|^{\alpha}+C\Big|\mathfrak{J}_{\beta}(t)|_{x=x_1}-\mathfrak{J}_{\beta}(t)|_{x=x_2}\Big|,
\end{align*}
and
\begin{align*}
&\Big|\mathfrak{J}_{\beta}(t)|_{x=x_1}-\mathfrak{J}_{\beta}(t)|_{x=x_2}\Big|
\\
&
\leq C\Big\{|x_2-x_1|\langle D_{x}w_1\rangle_{t,\Omega_{\tau}^{r}}^{(\frac{1+\alpha}{2}\nu)}\int\limits_{0}^{t}\frac{|w_2(x_2,s)-w_2(x_2,t)|s^{\frac{1+\alpha}{2}\nu}}{(t-s)^{1+\beta}}ds
\\&\quad+\int\limits_{0}^{t}\frac{|w_2(x_1,t)-w_2(x_2,t)+w_2(x_2,s)-w_{2}(x_1,s)|s^{\nu}\|\mathbf{D}_{t}^{\nu}w_1\|_{\C(\bar{\Omega}_{\tau}^{r})}}{(t-s)^{1+\beta}}ds
\Big\}\\
&
\leq C |x_1-x_2|^{\alpha}\tau^{\delta-\beta+\frac{\nu(1+\alpha)}{2}}(1+r^{1-\alpha})
\|w_2\|_{\C^{\delta}([0,T],\C^{1}(\bar{\Omega}))}\|w_1\|_{\C^{2+\alpha,\frac{2+\alpha}{2}\nu}(\bar{\Omega}_{\tau}^{r})},\end{align*}
which in turn entail
\begin{align}\label{4.3}\notag
&\|\xi\mathfrak{J}_{\beta}(t)\|_{\C([0,\tau],\C^{\alpha}(\bar{\Omega}^{r}))}
\\&\leq
C\tau^{\delta-\beta+\frac{\nu(1+\alpha)}{2}}
(r^{1-\alpha}+\tau^{(1-\alpha)\nu/2}r^{-\alpha}+1)\|w_1\|_{\C^{2+\alpha,\frac{2+\alpha}{2}\nu}(\bar{\Omega}^{r}_{\tau})}\|w_{2}\|_{\C^{\gamma}([0,T],\C^{1}(\bar{\Omega}))}.
\end{align}
Thus, taking into account Definition \ref{d2.1}, we will complete the proof of the first bound in (iii) if we obtain the corresponding estimate of the
seminorm $\langle \xi\mathfrak{J}_{\beta}(t)\rangle_{t,\Omega_{\tau}^{r}}^{(\alpha\nu/2)}$. To this end, assuming $T\geq t_{2}>t_1\geq 0$,
let us define
$$\Delta t=t_{2}-t_1.$$
If $\Delta t\geq t_1/2$, it is apparent that
$$
|\mathfrak{J}_{\beta}(t_2)-\mathfrak{J}_{\beta}(t_1)|\leq C\sum_{j=1}^{3}b_{j},
$$
where
\begin{align*}
b_{1}&=\int_{0}^{t_1}\frac{|w_2(x,t_1)-w_2(x,s)||w_1(x,s)-w_1(x,0)|}{(t_1-s)^{1+\beta}}ds,\\
b_{2}&=\int_{0}^{t_1}\frac{|w_2(x,t_2)-w_2(x,s)||w_1(x,s)-w_1(x,0)|}{(t_2-s)^{1+\beta}}ds,\\
b_{3}&=\int^{t_2}_{t_1}\frac{|w_2(x,t_2)-w_2(x,s)||w_1(x,s)-w_1(x,0)|}{(t_2-s)^{1+\beta}}ds.
\end{align*}
Then, appealing to the smoothness of the functions $w_1$ and $w_2$,
we conclude that
\begin{align*}
b_1&\leq C\int_{0}^{t_1}(t_1-s)^{\delta-\beta-1}s^{\nu}ds\langle w_2\rangle_{t,\Omega_T}^{(\delta)}\|\mathbf{D}^{\nu}_{t}w_1\|_{\C(\bar{\Omega}_{\tau}^{r})}\\
&
\leq C(\Delta t)^{\alpha\nu/2}\tau^{\delta+\nu-\beta-\alpha\nu/2} \|\mathbf{D}^{\nu}_{t}w_1\|_{\C(\bar{\Omega}_{\tau}^{r})}
\langle w_2\rangle_{t,\Omega_T}^{(\delta)},\end{align*}
and
\begin{align*}
b_3&\leq
\|\mathbf{D}^{\nu}_{t}w_1\|_{\C(\bar{\Omega}_{\tau}^{r})}
\langle w_2\rangle_{t,\Omega_T}^{(\delta)} t_{2}^{\nu}\int_{t_1}^{t_2}(t_2-s)^{\delta-1-\beta}ds\\
&
\leq C\tau^{\delta+\nu-\beta-\alpha\nu/2}(\Delta t)^{\alpha\nu/2}
\|\mathbf{D}^{\nu}_{t}w_1\|_{\C(\bar{\Omega}_{\tau}^{r})}
\langle w_2\rangle_{t,\Omega_T}^{(\gamma^{\star})}.
\end{align*}
The estimate for $b_2$ is analogous to the one of $b_1$.
Taking into account \eqref{4.3}, this yields the desired bound in (iii)
when $\Delta t\geq t_1/2$.
If instead $\Delta t<t_1/2$,
we rewrite the difference as
 \begin{equation}\label{4.5}
|\mathfrak{J}_{\beta}(t_2)-\mathfrak{J}_{\beta}(t_1)|\leq C\sum_{j=1}^{4}a_{j},
\end{equation}
where
\begin{align*}
a_{1}&=\int_{t_1-2\Delta t}^{t_1}\frac{|w_2(x,t_2)-w_2(x,s)||w_1(x,s)-w_1(x,0)|}{(t_2-s)^{1+\beta}}ds,\\
a_{2}&=\int_{t_1-2\Delta t}^{t_1}\frac{|w_2(x,t_1)-w_2(x,s)||w_1(x,s)-w_1(x,0)|}{(t_1-s)^{1+\beta}}ds,\\
a_{3}&=\int_{0}^{t_1-2\Delta t}|w_2(x,t_2)-w_2(x,s)||w_1(x,s)-w_1(x,0)|\Big|(t_2-s)^{-1-\beta}-(t_1-s)^{-1-\beta}\Big|ds,\\
a_{4}&=\int^{t_1-2\Delta t}_{0}\frac{|w_2(x,t_2)-w_2(x,t_1)||w_1(x,s)-w_1(x,0)|}{(t_1-s)^{1+\beta}}ds.
\end{align*}
Due to the properties of the functions $w_2$ and $w_1$, and exploiting  the mean-value theorem in the evaluation of  the term $a_3$, we end up with
\[
|\mathfrak{J}_{\beta}(t_2)-\mathfrak{J}_{\beta}(t_1)|\leq C\langle w_2\rangle_{t,\Omega_{T}}^{(\delta)} \|\mathbf{D}_{t}^{\nu}w_1\|_{\C(\bar{\Omega}_{\tau}^{r})}(\Delta t)^{\alpha\nu/2}\tau^{\delta+\nu-\beta-\alpha\nu/2},
\]
which completes the argument.
\end{proof}

\begin{thebibliography}{99}
\bibitem[KPSV1]{KPSV1} M. Krasnoschok, S. Pereverzyev, S.V. Siryk, N. Vasylyeva, Regularized reconstruction of the order in semilinear subdiffusion with memory, (In: Cheng J., Lu S., Yamamoto M. (Eds.) Inverse Problems and Related Topics ICIP2 2018), Springer Proceedings in Mathematics$\&$Statistics, \textbf{310} (2020) 205--236, doi:10.1007/978-981-15-1592-7-10.
\end{thebibliography}
\end{document}
